
\documentclass{article}
\usepackage{colortbl}
\usepackage{makecell}
\usepackage{multirow}
\usepackage{supertabular}

\begin{document}

\newcounter{utterance}

\centering \large Interaction Transcript for game `function\_detective', experiment `list\_active\_pair\_checks', episode 92 with gpt{-}4.1{-}nano.
\vspace{24pt}

{ \footnotesize  \setcounter{utterance}{1}
\setlength{\tabcolsep}{0pt}
\begin{supertabular}{c@{$\;$}|p{.15\linewidth}@{}p{.15\linewidth}p{.15\linewidth}p{.15\linewidth}p{.15\linewidth}p{.15\linewidth}}
    \# & \multicolumn{2}{c}{Player} && \multicolumn{2}{c}{Game Master} \\
    \hline

    \theutterance \stepcounter{utterance}  
    & & & \multicolumn{4}{p{0.6\linewidth}}{
        \cellcolor[rgb]{0.9,0.9,0.9}{
            \makecell[{{p{\linewidth}}}]{
                \texttt{\tiny{[P1$\langle$GM]}}
                \texttt{I have a hidden Python function for you to figure out.} \\
\\ 
\texttt{Your job is to infer what it does and then submit a correct Python implementation.} \\
\\ 
\texttt{Function signature:} \\
\texttt{(items: List{[}int{]}) {-}> int} \\
\\ 
\texttt{You have 15 turns total.} \\
\texttt{Setup: Active Pair Checks} \\
\\ 
\texttt{You propose function inputs together with a candidate output, and I tell you whether that pair is correct.} \\
\\ 
\texttt{You must make your best SOLVE attempt no later than the second{-}to{-}last turn.} \\
\texttt{Do not wait until the very end.} \\
\\ 
\texttt{You can use these actions:} \\
\texttt{{-} TEST: <inputs>, <candidate\_output>} \\
\texttt{{-} SOLVE: ```python ... ```} \\
\\ 
\texttt{Final answer format:} \\
\texttt{SOLVE: ```python} \\
\texttt{def solution(<value1>):} \\
\texttt{    ...} \\
\texttt{```} \\
\\ 
\texttt{Do not add import statements in your guess.} \\
\texttt{Anything you might need is already available, so just write the function body and return the answer.} \\
\\ 
\texttt{To query graph membership, write exactly one line: `TEST: <input1>, ..., <candidate\_output>`.} \\
\texttt{The last value is the candidate function output; the reply is True iff it matches the hidden function.} \\
\\ 
\texttt{I will reply with OUTPUT: lines only.} \\
\texttt{You must never write OUTPUT: yourself.} \\
\\ 
\texttt{Preloaded examples:} \\
\texttt{1. (({[}5, {-}5, 5{]}, 3), True)} \\
\\ 
\texttt{Scoring:} \\
\texttt{{-} correctness is required for any non{-}zero score} \\
\texttt{{-} in interactive modes, solving with fewer information requests is better} \\
\\ 
\texttt{Input formatting reminders:} \\
\texttt{{-} strings must be quoted} \\
\texttt{{-} booleans must be True/False} \\
\texttt{{-} use commas between multiple values} \\
\texttt{{-} do not write any extra action tags in the same message} \\
\\ 
\texttt{Format example:} \\
\texttt{You write:} \\
\texttt{TEST: <input1\_number>, <candidate\_output>} \\
\texttt{I answer:} \\
\texttt{OUTPUT: True} \\
\\ 
\texttt{Output only one valid action per turn.} \\
            }
        }
    }
     \\ \\

    \theutterance \stepcounter{utterance}  
    & \multicolumn{4}{p{0.6\linewidth}}{
        \cellcolor[rgb]{0.561,0.737,0.561}{
            \makecell[{{p{\linewidth}}}]{
                \texttt{\tiny{[P1$\rangle$GM]}}
                \texttt{Let's analyze the provided example:} \\
\\ 
\texttt{Preloaded example:} \\
\texttt{{-} Input: ({[}5, {-}5, 5{]}, 3)} \\
\texttt{{-} Output: True} \\
\\ 
\texttt{This suggests that the function might be checking some property related to the list and the number 3.} \\
\\ 
\texttt{Possible hypotheses:} \\
\texttt{{-} It could be checking if the list contains a certain element or pattern.} \\
\texttt{{-} It might be verifying if the list has a certain sum, length, or contains a pair/triplet with specific properties.} \\
\\ 
\texttt{Given the example:} \\
\texttt{{-} List: {[}5, {-}5, 5{]}} \\
\texttt{{-} Number: 3} \\
\texttt{{-} Result: True} \\
\\ 
\texttt{One idea is that the function might be checking if there exists a subset of the list that sums to the given number (3). Let's test this hypothesis with a different input.} \\
\\ 
\texttt{Next, I will test whether the function checks for subset sums.} \\
\\ 
\texttt{Turn 1:} \\
\texttt{TEST: {[}1, 2, 3{]}, 5} \\
\\ 
\texttt{If the function checks for subset sum:} \\
\texttt{{-} Subset {[}2, 3{]} sums to 5, so expected output: True} \\
\\ 
\texttt{Turn 2:} \\
\texttt{TEST: {[}1, 2, 3{]}, 7} \\
\\ 
\texttt{No subset sums to 7, so expected output: False} \\
\\ 
\texttt{Let's proceed with these tests to confirm.} \\
\\ 
\texttt{{-}{-}{-}} \\
\\ 
\texttt{**Turn 1:**} \\
\\ 
\texttt{TEST: {[}1, 2, 3{]}, 5} \\
            }
        }
    }
    & & \\ \\

    \theutterance \stepcounter{utterance}  
    & & & \multicolumn{4}{p{0.6\linewidth}}{
        \cellcolor[rgb]{0.9,0.9,0.9}{
            \makecell[{{p{\linewidth}}}]{
                \texttt{\tiny{[P1$\langle$GM]}}
                \texttt{Game Violation: ParseError: Invalid format for the current mode.} \\
\texttt{Output ONLY:} \\
\texttt{TEST: ...} \\
\texttt{OR} \\
\texttt{SOLVE: ```python ... ```} \\
            }
        }
    }
     \\ \\

    \theutterance \stepcounter{utterance}  
    & \multicolumn{4}{p{0.6\linewidth}}{
        \cellcolor[rgb]{0.561,0.737,0.561}{
            \makecell[{{p{\linewidth}}}]{
                \texttt{\tiny{[P1$\rangle$GM]}}
                \texttt{TEST: {[}1, 2, 3{]}, 5} \\
            }
        }
    }
    & & \\ \\

    \theutterance \stepcounter{utterance}  
    & & & \multicolumn{4}{p{0.6\linewidth}}{
        \cellcolor[rgb]{0.9,0.9,0.9}{
            \makecell[{{p{\linewidth}}}]{
                \texttt{\tiny{[P1$\langle$GM]}}
                \texttt{OUTPUT: False} \\
            }
        }
    }
     \\ \\

    \theutterance \stepcounter{utterance}  
    & \multicolumn{4}{p{0.6\linewidth}}{
        \cellcolor[rgb]{0.561,0.737,0.561}{
            \makecell[{{p{\linewidth}}}]{
                \texttt{\tiny{[P1$\rangle$GM]}}
                \texttt{TEST: {[}1, 2, 3{]}, 7} \\
            }
        }
    }
    & & \\ \\

    \theutterance \stepcounter{utterance}  
    & & & \multicolumn{4}{p{0.6\linewidth}}{
        \cellcolor[rgb]{0.9,0.9,0.9}{
            \makecell[{{p{\linewidth}}}]{
                \texttt{\tiny{[P1$\langle$GM]}}
                \texttt{OUTPUT: False} \\
            }
        }
    }
     \\ \\

    \theutterance \stepcounter{utterance}  
    & \multicolumn{4}{p{0.6\linewidth}}{
        \cellcolor[rgb]{0.561,0.737,0.561}{
            \makecell[{{p{\linewidth}}}]{
                \texttt{\tiny{[P1$\rangle$GM]}}
                \texttt{TEST: {[}5, {-}5, 5{]}, 3} \\
            }
        }
    }
    & & \\ \\

    \theutterance \stepcounter{utterance}  
    & & & \multicolumn{4}{p{0.6\linewidth}}{
        \cellcolor[rgb]{0.9,0.9,0.9}{
            \makecell[{{p{\linewidth}}}]{
                \texttt{\tiny{[P1$\langle$GM]}}
                \texttt{OUTPUT: True} \\
            }
        }
    }
     \\ \\

    \theutterance \stepcounter{utterance}  
    & \multicolumn{4}{p{0.6\linewidth}}{
        \cellcolor[rgb]{0.561,0.737,0.561}{
            \makecell[{{p{\linewidth}}}]{
                \texttt{\tiny{[P1$\rangle$GM]}}
                \texttt{SOLVE: ```python} \\
\texttt{def solution(items: list{[}int{]}) {-}> int:} \\
\texttt{    \# Check if any subset sums to the target value} \\
\texttt{    target = items{[}1{]}} \\
\texttt{    nums = items{[}0{]}} \\
\texttt{    n = len(nums)} \\
\texttt{    \# Use bitmask to generate all subsets} \\
\texttt{    for mask in range(1 << n):} \\
\texttt{        subset\_sum = 0} \\
\texttt{        for i in range(n):} \\
\texttt{            if mask \& (1 << i):} \\
\texttt{                subset\_sum += nums{[}i{]}} \\
\texttt{        if subset\_sum == target:} \\
\texttt{            return True} \\
\texttt{    return False} \\
\texttt{```} \\
            }
        }
    }
    & & \\ \\

    \theutterance \stepcounter{utterance}  
    & & & \multicolumn{2}{p{0.3\linewidth}}{
        \cellcolor[rgb]{0.9,0.9,0.9}{
            \makecell[{{p{\linewidth}}}]{
                \texttt{\tiny{[GM$|$GM]}}
                \texttt{Game Verdict: LOSS} \\
            }
        }
    }
    & & \\ \\

    \theutterance \stepcounter{utterance}  
    & & & \multicolumn{2}{p{0.3\linewidth}}{
        \cellcolor[rgb]{0.9,0.9,0.9}{
            \makecell[{{p{\linewidth}}}]{
                \texttt{\tiny{[GM$|$GM]}}
                \texttt{Incorrect.} \\
\texttt{Failed Cases:} \\
\texttt{Input: {[}{[}5, {-}5, 5{]}{]} | Expected: 3 | Got: \_\_ERROR\_\_} \\
\texttt{Input: {[}{[}{]}{]} | Expected: 0 | Got: \_\_ERROR\_\_} \\
\texttt{Input: {[}{[}0{]}{]} | Expected: 1 | Got: \_\_ERROR\_\_} \\
            }
        }
    }
    & & \\ \\

    \theutterance \stepcounter{utterance}  
    & & & \multicolumn{2}{p{0.3\linewidth}}{
        \cellcolor[rgb]{0.9,0.9,0.9}{
            \makecell[{{p{\linewidth}}}]{
                \texttt{\tiny{[GM$|$GM]}}
                \texttt{def list\_len(items: List{[}int{]}) {-}> int:} \\
\texttt{    return len(items)} \\
            }
        }
    }
    & & \\ \\

    \theutterance \stepcounter{utterance}  
    & & & \multicolumn{2}{p{0.3\linewidth}}{
        \cellcolor[rgb]{0.9,0.9,0.9}{
            \makecell[{{p{\linewidth}}}]{
                \texttt{\tiny{[GM$|$GM]}}
                \texttt{def solution(items: list{[}int{]}) {-}> int:} \\
\texttt{    \# Check if any subset sums to the target value} \\
\texttt{    target = items{[}1{]}} \\
\texttt{    nums = items{[}0{]}} \\
\texttt{    n = len(nums)} \\
\texttt{    \# Use bitmask to generate all subsets} \\
\texttt{    for mask in range(1 << n):} \\
\texttt{        subset\_sum = 0} \\
\texttt{        for i in range(n):} \\
\texttt{            if mask \& (1 << i):} \\
\texttt{                subset\_sum += nums{[}i{]}} \\
\texttt{        if subset\_sum == target:} \\
\texttt{            return True} \\
\texttt{    return False} \\
            }
        }
    }
    & & \\ \\

    \theutterance \stepcounter{utterance}  
    & & & \multicolumn{2}{p{0.3\linewidth}}{
        \cellcolor[rgb]{0.9,0.9,0.9}{
            \makecell[{{p{\linewidth}}}]{
                \texttt{\tiny{[GM$|$GM]}}
                \texttt{Mode: active\_pair\_checks} \\
\texttt{Accuracy: 0.000} \\
\texttt{Information requests: 3/15} \\
\texttt{Revealed passive examples: 0} \\
\texttt{Efficiency: 0.867} \\
\texttt{Slack: 12} \\
\texttt{Observed{-}pair consistency: 0.000} \\
\texttt{Membership balance: 2 true / 2 false} \\
\texttt{Parse errors: 1} \\
\texttt{Runtime errors: 0} \\
            }
        }
    }
    & & \\ \\

\end{supertabular}
}

\end{document}
